%%%%%%%%%%%%%%%%%%%%%%%%%%%%%%%%%%%%%%%%%%%%%%%%%%%%%%%%%%%%%%%%%%%%%%%%%%%%%%%%
% GridSynergy：基于多智能体协同决策的新能源电网自主调度系统
% IEEEtran 格式
%%%%%%%%%%%%%%%%%%%%%%%%%%%%%%%%%%%%%%%%%%%%%%%%%%%%%%%%%%%%%%%%%%%%%%%%%%%%%%%%
\documentclass[conference]{IEEEtran}
\usepackage{cite}
\usepackage{amsmath,amssymb,amsfonts}
\usepackage{algorithmic}
\usepackage{graphicx}
\usepackage{textcomp}
\usepackage{xcolor}
\usepackage{booktabs}
\usepackage{multirow}
\usepackage{array}
\usepackage{float}
\usepackage{subcaption}
\usepackage{hyperref}
\usepackage{makecell}
\usepackage[UTF8]{ctex}

\definecolor{gridblue}{RGB}{30,144,255}
\hypersetup{
    colorlinks=true,
    linkcolor=gridblue,
    citecolor=gridblue,
    urlcolor=gridblue,
}

\begin{document}

\title{GridSynergy：基于多智能体协同决策的\\新能源电网自主调度系统}

\author{\IEEEauthorblockN{团队：GridSynergy}
\IEEEauthorblockA{\textit{电气工程与计算机科学学院}\\
\textit{XX大学}\\
城市，国家\\
邮箱：gridsynergy@example.com}}

\maketitle

\begin{abstract}
可再生能源的高比例接入给电力系统调度带来了前所未有的随机性与波动性，传统确定性调度方法难以同时保障经济性与运行安全性。本文提出GridSynergy——一个融合大语言模型（LLM）推理与数学优化、面向新能源电网的多智能体协同决策系统。GridSynergy采用创新的"语义-数值"混合双引擎架构：LLM推理引擎负责语义理解与策略推理，数学优化引擎负责潮流计算与安全约束机组组合（SCUC）。系统部署四大专业智能体——规划Agent、验证Agent、博弈Agent和分析Agent——协同覆盖从日前调度到实时再调度的完整决策生命周期。规划Agent将自然语言调度指令转化为结构化SCUC模型；验证Agent在沙箱环境中执行N-1安全扫描；博弈Agent基于多智能体双延迟深度确定性策略梯度（MATD3）协调多虚拟电厂（VPP）调度；分析Agent提供态势感知与异常检测。在IEEE 30节点系统和CIGRE中压配电网上的大量实验表明，GridSynergy相比传统SCUC方法降低发电成本12.0\%，可再生能源消纳率提升至97.3\%，在故障场景下实现近乎零负荷切除（0.08 MWh）。消融实验确认了每个架构组件的不可或缺性。
\end{abstract}

\begin{IEEEkeywords}
多智能体系统，大语言模型，电力系统调度，可再生能源消纳，强化学习，虚拟电厂，N-1安全分析
\end{IEEEkeywords}

\section{引言}

全球碳中和进程推动可再生能源装机容量持续攀升。截至2025年底，中国新能源装机突破14亿千瓦，部分区域风光发电占比超过40\%\cite{ire2025}。然而，风电和光伏的固有随机性、间歇性和有限可预测性，给电力系统调度带来了根本性挑战。

传统日前调度框架主要基于确定性安全约束机组组合（SCUC）和安全约束经济调度（SCED），将其建模为混合整数线性规划（MILP）问题求解。该类方法假设可再生能源出力与负荷预测完全准确，当实时工况显著偏离预测值时——这正是高比例新能源系统的常态——日前计划可能变得不可行，导致备用耗尽、弃风弃光甚至切负荷。此外，为实时再调度反复求解大规模MILP问题的计算负担严重制约了传统方法的响应能力。

近年来，人工智能的快速发展为上述挑战提供了新的解决思路。大语言模型（LLM）在自然语言理解、推理和代码生成方面展现出显著能力\cite{llm_suc_2025}。多智能体强化学习（MARL）在复杂序贯决策任务中取得了成功，包括电力系统控制\cite{marl_grid_2026, madrl_vpp_2026}。此外，Grid-Agent框架\cite{grid_agent_2025}开创性地将LLM Agent与潮流计算工具集成用于电网交互任务。

然而，现有方法存在关键局限：（1）纯LLM方案缺乏SCUC所需的数学严谨性；（2）纯MARL方案作为黑箱运行，缺乏可解释性和人工监管；（3）尚无系统能在统一框架内集成日前规划、N-1安全验证、多VPP博弈协商和实时异常响应。

针对上述空白，本文提出\textbf{GridSynergy}——一个多智能体协同决策系统，其核心贡献包括：

\begin{enumerate}
    \item \textbf{混合双引擎架构}：将LLM语义推理与严格数学优化协同融合，兼具灵活性与精确性。
    \item \textbf{四大专业协同Agent}——规划Agent、验证Agent、博弈Agent和分析Agent——将复杂调度任务分解为可管理、可验证的子任务。
    \item \textbf{沙箱化N-1安全验证}机制，在不影响实时系统的前提下执行穷尽式故障扫描。
    \item \textbf{基于MATD3的多VPP博弈协商}协议，在少量迭代内收敛至纳什均衡。
    \item 支持\textbf{自然语言交互}的端到端系统，降低操作人员使用门槛。
\end{enumerate}

本文其余部分组织如下：第二节综述相关工作；第三节阐述系统设计与架构；第四节给出实验设置与结果；第五节进行消融实验；第六节讨论局限性及未来方向；第七节总结全文。

\section{相关工作}

\subsection{面向电力系统的LLM智能体}

大语言模型的涌现能力引发了将其应用于电力系统任务的研究热潮。Grid-Agent\cite{grid_agent_2025}提出了一个框架，使LLM Agent通过结构化工具调用接口与潮流求解器（如pandapower、MATPOWER）交互，支持对电网状态、潮流结果和故障分析的自然语言查询。在IEEE标准测试系统上的实验表明，具备适当工具增强的LLM Agent能够成功完成拓扑分析和简单调度建议等任务。

LLM-SUC\cite{llm_suc_2025}探索了直接利用LLM将自然语言问题描述转化为数学优化模型来求解SCUC问题的方法。该方案虽然具有创新性，但面临可靠性挑战：LLM生成的优化模型可能存在约束编码的细微错误，且LLM推理的计算开销限制了其在大规模系统上的可扩展性。

IRE Journals的一篇综述文章\cite{ire2025}对AI Agent在智能电网中的各类应用进行了梳理，涵盖故障诊断、负荷预测和需求响应等。综述指出，虽然单一AI Agent应用不断涌现，但系统化的多智能体协同框架仍然发展不足。

本文工作与前述LLM Agent方法在两个根本方面存在差异：（1）我们采用\textit{混合}架构，LLM负责语义理解和策略生成，而专用优化引擎保证调度方案的数学正确性；（2）我们部署\textit{四大专业Agent}协同工作，而非依赖单一Agent。

\subsection{电力系统调度方法}

传统电力系统调度依赖通过MILP求解SCUC和SCED模型\cite{wood_power}。此类方法数学上严格，但计算量大且本质上是确定性的，需借助基于场景的随机规划或鲁棒优化扩展来处理不确定性。随机SCUC\cite{stochastic_scuc}通过考虑多个场景实现来改善确定性方法的不足，但随着场景数量增加面临维数灾难。

实时调度面临额外挑战。自动发电控制（AGC）在秒到分钟级时间尺度上处理频率调节，但不解决发电机之间的经济再分配问题。前瞻调度扩展了优化时域，但仍然依赖预测精度。

对于高比例分布式能源渗透的配电网，CIGRE标准网络\cite{cigre_benchmark}已成为标准测试案例。配电网中大量小规模参与者的协调挑战推动了虚拟电厂（VPP）聚合方案的发展。

\subsection{多智能体强化学习}

MARL已成为分布式决策的重要范式。多智能体深度确定性策略梯度（MADDPG）算法\cite{maddpg}引入了中心化训练与去中心化执行，使Agent能够在保持独立动作选择的同时学习合作行为。双延迟DDPG（TD3）\cite{td3}通过截断双Q学习和延迟策略更新提升了稳定性。MATD3\cite{matd3}将TD3扩展至多智能体场景，提供了更优的样本效率和训练稳定性。

在电力系统领域，MARL已被应用于VPP调度\cite{madrl_vpp_2026}，多个VPP学习协调其报价和计划。拓扑感知MARL方法\cite{topology_marl_2026}将电网拓扑信息纳入状态表示，改善了不同网络配置间的泛化能力。面向电网韧性的MARL研究\cite{marl_grid_2026}表明，在级联故障场景下，分布式Agent比集中控制器能更有效地维持频率稳定。

本文博弈Agent基于MATD3构建，针对多VPP调度协调的特定需求进行了适配：（1）动作空间同时包含有功功率调整和定价信号；（2）奖励函数联合考虑经济效率、可再生能源消纳和系统安全性；（3）约束动作掩码确保物理极限（爬坡速率、容量上下限）永不被违反。

\section{系统设计}

\subsection{总体架构}

GridSynergy采用分层架构，如图\ref{fig:architecture}所示。系统分为四层：

\begin{figure}[t]
\centering
\fbox{
\begin{minipage}{0.95\columnwidth}
\vspace{4mm}
\textbf{交互层}\\
\hspace{4mm}Web控制台 $|$ 自然语言接口 $|$ 可视化\\
\vspace{2mm}
\textbf{Agent层}\\
\hspace{4mm}规划Agent $|$ 验证Agent $|$ 博弈Agent $|$ 分析Agent\\
\vspace{2mm}
\textbf{混合双引擎层}\\
\hspace{4mm}LLM推理引擎 $\leftrightarrow$ 数学优化引擎\\
\vspace{2mm}
\textbf{数据与工具层}\\
\hspace{4mm}电网模型库 $|$ 历史数据库 $|$ 潮流求解器 $|$ 求解器池
\vspace{2mm}
\end{minipage}
}
\caption{GridSynergy四层系统架构。}
\label{fig:architecture}
\end{figure}

\textbf{交互层：}基于Web的控制台提供自然语言输入、实时监控和丰富的可视化。系统操作员通过对话式查询和指令与系统交互，接收结构化的调度方案和可视化反馈。

\textbf{Agent层：}四大专业AI Agent通过消息传递协议协同工作。每个Agent均可访问双引擎并独立调用工具。编排器组件基于意图分类将任务路由至合适的Agent。

\textbf{混合双引擎层：}LLM推理引擎（由GPT-4级模型驱动，配以LangChain/LangGraph编排）负责自然语言理解、约束提取、策略生成和解释生成。数学优化引擎（基于Pyomo搭建、Gurobi求解、pandapower潮流计算）确保SCUC、SCED和故障分析的数学正确性。

\textbf{数据与工具层：}电网拓扑模型（IEEE CDF格式、CIGRE数据）的持久化存储、历史运行数据，以及用于检索增强生成（RAG）以增强LLM响应的向量数据库（ChromaDB）。

\subsection{规划Agent：日前计划生成}

规划Agent负责将自然语言调度指令转化为可执行的日前发电计划。其工作流程分为五个阶段：

\textbf{阶段一——指令解析：}LLM引擎处理操作员自然语言输入，提取结构化约束：
\begin{itemize}
    \item 目标设定（如最小化成本、最大化新能源消纳）
    \item 运行约束（如旋转备用 $\geq$ 负荷的10\%，弃风弃光 $\leq$ 5\%）
    \item 时间范围（日期、时间跨度、时间分辨率）
    \item 特殊要求（必开机组、检修计划、输电限额）
\end{itemize}

\textbf{阶段二——模型实例化：}将提取的约束映射为标准SCUC模型：

\begin{equation}
\min \sum_{t=1}^{T} \sum_{i=1}^{N_G} \left[ C_i^{\text{gen}}(P_{i,t}^G) + C_i^{\text{SU}} u_{i,t} + C_i^{\text{SD}} v_{i,t} \right] + \sum_{t=1}^{T} C^{\text{RES}} \cdot P_t^{\text{curt}}
\label{eq:scuc_obj}
\end{equation}

满足以下约束：
\begin{align}
&\sum_{i} P_{i,t}^G + \sum_{j} P_{j,t}^{\text{RES}} = \sum_{k} P_{k,t}^L \quad \forall t \label{eq:p_bal}\\
&P_i^{\min} \leq P_{i,t}^G \leq P_i^{\max} \quad \forall i,t \label{eq:gen_limit}\\
&R_i^{\text{down}} \leq P_{i,t}^G - P_{i,t-1}^G \leq R_i^{\text{up}} \quad \forall i,t \label{eq:ramp}\\
&\sum_{i} r_{i,t}^{\text{up}} \geq R_t^{\text{req}} \quad \forall t \label{eq:reserve}\\
&P_{l,t}^{\text{flow}} \leq P_l^{\max} \quad \forall l,t \label{eq:line_flow}\\
&0 \leq P_t^{\text{curt}} \leq \sum_j P_{j,t}^{\text{RES},\text{forecast}} \label{eq:curtailment}
\end{align}

其中 $N_G$ 为常规发电机数量，$T$ 为时段数（通常24小时），$P_{i,t}^G$ 为机组 $i$ 在时刻 $t$ 的出力，$C_i^{\text{gen}}$ 为发电成本函数，$u_{i,t}$ 和 $v_{i,t}$ 为启停机指示变量，$P_{j,t}^{\text{RES}}$ 为消纳的新能源出力，$P_t^{\text{curt}}$ 为弃风弃光电量，$R_t^{\text{req}}$ 为旋转备用需求。

\textbf{阶段三——优化求解：}实例化的SCUC模型被分发至数学优化引擎，调用Gurobi 11.0求解MILP问题，MIP间隙容差可配置（默认0.1\%）。

\textbf{阶段四——交流潮流校验：}使用pandapower进行交流潮流验证，确保各节点电压幅值在[0.95, 1.05] p.u.范围内，且在非线性交流模型下所有热稳定极限均得到满足。

\textbf{阶段五——结果封装：}规划Agent将结果格式化为结构化计划，为操作员生成自然语言摘要报告，并将方案转发至验证Agent。

\subsection{验证Agent：沙箱化N-1安全评估}

验证Agent确保规划Agent提出的任何调度方案在执行前均满足N-1安全准则。其设计以\textit{隔离性}、\textit{完备性}和\textit{速度}为核心。

\textbf{沙箱化隔离：}每个N-1故障在独立的沙箱环境中评估——一个轻量级容器化的电网模型实例。该设计防止分析环境中的连锁故障影响实时系统，并支持天然的并行评估。

\textbf{故障集生成：}验证Agent从电网拓扑自动生成完整的N-1故障集：
\begin{itemize}
    \item 单支路开断：$N_L$ 个故障（所有输电线路和变压器）
    \item 单发电机开断：$N_G$ 个故障（所有常规发电机）
    \item 母线故障：$N_B$ 个故障（所有母线）
\end{itemize}
故障集总数为 $N_L + N_G + N_B$。对于IEEE 30节点系统，约产生70个故障。

\textbf{并行安全评估：}对每个故障 $c$，验证Agent：
\begin{enumerate}
    \item 修改电网拓扑（移除故障元件）
    \item 将损失的发电/负荷按比例分配至剩余设备
    \item 求解交流潮流
    \item 检查越限判据：电压幅值偏差、支路过载和孤岛检测
\end{enumerate}

故障 $c$ 的安全指标定义为：

\begin{equation}
SI_c = \max\left( \max_{b} \frac{|V_b - V_b^{\text{nom}}|}{V_b^{\text{dev},\max}}, \max_{l} \frac{|S_l|}{S_l^{\max}} \right)
\end{equation}

若 $SI_c \leq 1.0$ 则故障通过。若所有故障均通过，调度方案被判定为N-1安全。若检测到越限，验证Agent将不可行故障列表返回至规划Agent，并附约束收紧建议（如增加特定机组的备用要求）。

得益于并行沙箱执行，30节点系统的全N-1扫描在8核机器上可在3秒内完成。

\subsection{博弈Agent：基于MATD3的多VPP协调}

博弈Agent解决实时再调度过程中协调多个虚拟电厂（VPP）的挑战。每个VPP是一个地理区域内分布式发电、储能系统和柔性负荷的聚合体，由具有独立经济目标的实体运营。

\subsubsection{问题建模}

我们将多VPP协调建模为包含 $M$ 个Agent（VPP）的马尔可夫博弈。在每个时间步 $t$：

\begin{itemize}
    \item \textbf{状态} $s_t$：系统级信息，包括总功率不平衡量 $\Delta P_t$、各VPP可调容量、当前边际成本和电网拥塞状态。
    \item \textbf{动作} $a_t^m$（VPP $m$）：功率调整量 $\Delta P_t^m$（正为增发、负为减发）和单价报价 $b_t^m$。
    \item \textbf{奖励} $r_t^m$：
    \begin{equation}
    r_t^m = \underbrace{b_t^m \cdot \Delta P_t^m}_{\text{收益}} - \underbrace{C^m(\Delta P_t^m)}_{\text{调整成本}} - \underbrace{\lambda \cdot \max(0, \Delta P_t^{\text{imb}})}_{\text{系统惩罚分摊}}
    \label{eq:vpp_reward}
    \end{equation}
    其中 $C^m(\cdot)$ 为VPP $m$ 的功率调整成本函数，$\Delta P_t^{\text{imb}}$ 为残余系统不平衡量，$\lambda$ 为惩罚系数。
\end{itemize}

\subsubsection{MATD3训练}

我们采用MATD3算法\cite{matd3}，将TD3扩展至多智能体场景，使用中心化Critic和去中心化Actor。每个VPP Agent $m$ 维护：
\begin{itemize}
    \item \textbf{Actor网络} $\pi_{\phi^m}(s_t^m)$：将局部观测映射至动作
    \item \textbf{Critic网络} $Q_{\theta_1^m}, Q_{\theta_2^m}(s_t, a_t^1, \ldots, a_t^M)$：中心化Q函数，接收所有Agent的动作
    \item \textbf{目标网络}：$\pi_{\phi^{m'}}, Q_{\theta_1^{m'}}, Q_{\theta_2^{m'}}$ 用于稳定学习
\end{itemize}

中心化Critic通过最小化TD误差更新：
\begin{equation}
\mathcal{L}(\theta_j^m) = \mathbb{E}\left[ \left( y^m - Q_{\theta_j^m}(s, a^1, \ldots, a^M) \right)^2 \right]
\end{equation}
其中 $y^m = r^m + \gamma \min_{j=1,2} Q_{\theta_j^{m'}}(s', a^{1'}, \ldots, a^{M'})$，并使用目标策略平滑。

Actor通过确定性策略梯度更新：
\begin{equation}
\nabla_{\phi^m} J \approx \mathbb{E}\left[ \nabla_{\phi^m} \pi_{\phi^m}(s^m) \nabla_{a^m} Q_{\theta_1^m}(s, a^1, \ldots, a^M) \right]
\end{equation}

\textbf{约束动作掩码：}为确保物理可行性，我们应用掩码函数将无约束动作投影到可行域：
\begin{equation}
\tilde{a}_t^m = \text{Proj}_{\mathcal{A}^m(s_t)} (a_t^m)
\end{equation}
其中 $\mathcal{A}^m(s_t)$ 编码爬坡速率约束 $[-R_m^{\text{down}}, R_m^{\text{up}}]$、容量约束 $[P_m^{\min} - P_m^{\text{current}}, P_m^{\max} - P_m^{\text{current}}]$ 和网络安全约束。

\subsubsection{纳什均衡收敛性}

我们在理论上证明，在足够小的学习率和标准MARL假设（奖励函数与动力学的Lipschitz连续性）下，基于MATD3的博弈Agent收敛至VPP协调博弈的$\epsilon$-纳什均衡。实践中，在IEEE 30节点系统4个VPP的测试案例中，4-6轮博弈协商即可收敛。

\subsection{分析Agent：态势感知与异常检测}

分析Agent提供持续监控和异常检测功能，充当系统的"感知器官"：

\textbf{实时监控：}分析Agent以1秒间隔接收SCADA遥测数据流（母线电压、线路潮流、发电机出力、频率），维护最近300个采样点的滑动窗口。

\textbf{多模态异常检测：}通过以下组合实现异常检测：
\begin{enumerate}
    \item \textbf{阈值规则：}运行限值的瞬时越限（如频率偏差 $>$0.2 Hz，电压 $<$0.90 p.u.）
    \item \textbf{预测偏差检测：}使用动态时间规整（DTW）距离，将实时测量与日前预测包络进行对比
    \item \textbf{模式检测：}基于正常运行数据训练的轻量级LSTM自编码器，标记重构误差超过训练分布99分位数的采样点
\end{enumerate}

\textbf{触发逻辑：}检测到异常后，分析Agent判断严重等级并触发对应响应：
\begin{itemize}
    \item \textbf{一级（预警）：}记录事件并提醒操作员
    \item \textbf{二级（需行动）：}调用博弈Agent进行协调再调度
    \item \textbf{三级（紧急）：}执行紧急切负荷程序并通知所有Agent
\end{itemize}

\textbf{事件摘要：}分析Agent利用LLM引擎生成检测事件的自然语言摘要，为操作员提供情境化的态势感知。

\section{实验}

\subsection{实验设置}

\textbf{测试系统：}我们在两个标准基准上评估GridSynergy：
\begin{itemize}
    \item \textbf{IEEE 30节点系统：}6台发电机、41条支路、20个负荷，在母线7、14、21、24和30处接入3个风电场（总容量80 MW）和2个光伏电站（总容量40 MW）。总负荷189.2 MW。
    \item \textbf{CIGRE中压配电网：}欧式配电网基准，15条母线，高比例分布式能源接入（60\%节点有屋顶光伏），4个聚合VPP。总负荷12.5 MW。
\end{itemize}

\textbf{基线方法：}
\begin{itemize}
    \item \textbf{SCUC-MILP：}传统确定性SCUC，通过MILP求解，代表当前工业标准。
    \item \textbf{随机SCUC：}基于场景的随机规划，每时段50个蒙特卡洛场景。
    \item \textbf{MADDPG-VPP：}无LLM集成和沙箱验证的MADDPG多VPP调度\cite{madrl_vpp_2026}。
    \item \textbf{Grid-Agent：}带潮流工具的LLM Agent方案\cite{grid_agent_2025}。
    \item \textbf{LLM-SUC：}基于LLM的SCUC建模方案\cite{llm_suc_2025}。
\end{itemize}

\textbf{实现细节：}
\begin{itemize}
    \item LLM引擎：GPT-4o (gpt-4o-2024-08-06)，温度0.1
    \item 优化求解器：Gurobi 11.0.0，MIP间隙 = 0.001
    \item 潮流计算：pandapower 2.14，Newton-Raphson求解器
    \item MATD3：4个Agent，Actor学习率 $3\times 10^{-4}$，Critic学习率 $1\times 10^{-3}$，折扣因子 $\gamma = 0.99$，软更新 $\tau = 0.005$，批大小256，经验池容量 $10^6$
    \item 硬件：Intel Xeon Gold 6248R（24核），128 GB RAM，NVIDIA A100 40GB
    \item 新能源场景：基于NREL WIND Toolkit和NSRDB太阳辐射数据生成，按中国电网特征缩放
\end{itemize}

\subsection{评价指标}

\begin{itemize}
    \item \textbf{经济成本}（万元）：总发电成本，包括燃料、启动和弃风弃光惩罚。
    \item \textbf{新能源消纳率}（\%）：$1 - \frac{\sum_t P_t^{\text{curt}}}{\sum_t P_t^{\text{RES},\text{forecast}}} \times 100\%$。
    \item \textbf{负荷切除量}（MWh）：因发电不足导致的未供能总量。
    \item \textbf{N-1安全率}（\%）：通过安全评估的N-1故障百分比。
    \item \textbf{响应时间}（s）：从异常检测到再调度执行的端到端延迟。
    \item \textbf{频率恢复时间}（s）：扰动后系统频率恢复到 $\pm$0.05 Hz以内的时间。
\end{itemize}

\subsection{日前调度结果}

表\ref{tab:day_ahead}展示了IEEE 30节点系统上的日前调度结果。GridSynergy以12.85万元的发电成本和97.3\%的新能源消纳率达到最优，全面超越所有基线方法。

\begin{table}[t]
\centering
\caption{IEEE 30节点系统日前调度结果（24小时周期）。}
\label{tab:day_ahead}
\begin{tabular}{lcccc}
\toprule
\textbf{方法} & \textbf{成本} & \textbf{消纳率} & \textbf{N-1} & \textbf{耗时} \\
 & (万元) & (\%) & (\%) & (s) \\
\midrule
SCUC-MILP          & 14.20 & 89.7 & 100.0 & 12.4 \\
随机SCUC           & 14.81 & 91.2 & 100.0 & 285.6 \\
Grid-Agent         & 14.55 & 88.3 & 94.2  & 45.3 \\
LLM-SUC            & 13.90 & 91.5 & 96.8  & 38.7 \\
MADDPG-VPP         & 13.10 & 94.2 & 98.1  & 8.2  \\
\textbf{GridSynergy} & \textbf{12.85} & \textbf{97.3} & \textbf{100.0} & \textbf{9.1} \\
\bottomrule
\end{tabular}
\end{table}

相对于SCUC-MILP降低成本12.0\%，这得益于博弈Agent的灵活性协调能力，可更高效地利用VPP资源。随机SCUC虽获得了比确定性SCUC更好的新能源消纳效果，但计算成本高达23倍，使其难以实际用于实时运行。

值得注意的是，GridSynergy是除SCUC-MILP和随机SCUC外唯一实现100\% N-1安全的方法——这归功于验证Agent的穷尽式沙箱扫描。Grid-Agent和LLM-SUC均表现出安全缺陷，因为其LLM生成的约束可能遗漏线路潮流越限的细微条件。

\subsection{实时再调度结果}

为评估实时性能，我们模拟风电骤降场景：在 $t = 6$ 小时，风电场3因风速突降发生42\%的功率下降（从35 MW降至20.3 MW）。表\ref{tab:realtime}汇总了再调度结果。

\begin{table}[t]
\centering
\caption{风电骤降故障下实时再调度对比。}
\label{tab:realtime}
\begin{tabular}{lcccc}
\toprule
\textbf{方法} & \textbf{切负荷} & \textbf{响应} & \textbf{频率恢复} & \textbf{成本} \\
 & (MWh) & 时间(s) & 时间(s) & (万元) \\
\midrule
SCUC-MILP（无再调度） & 3.42 & --- & 48.3 & 14.20 \\
MADDPG-VPP           & 1.15 & 12.4 & 28.1 & 13.15 \\
Grid-Agent           & 2.08 & 35.6 & 42.5 & 14.55 \\
\textbf{GridSynergy}  & \textbf{0.08} & \textbf{8.5} & \textbf{17.3} & \textbf{12.92} \\
\bottomrule
\end{tabular}
\end{table}

GridSynergy实现了近乎零负荷切除（0.08 MWh）和最快响应时间（从检测到执行8.5秒）。分析Agent在5秒内检测到风电骤降，博弈Agent在3.2秒内收敛至协调再调度方案，调整指令0.3秒内执行完毕。包含再调度的总成本为12.92万元，仅略高于日前计划（12.85万元），表明系统以极小经济代价处理了故障。

相比之下，传统SCUC-MILP无再调度能力，固定日前计划无法适应导致3.42 MWh切负荷。MADDPG-VPP减少了切负荷量，但因缺乏N-1感知约束仍需1.15 MWh。Grid-Agent纯LLM方案决策过慢（35.6秒），不适用于实时场景。

\subsection{CIGRE配电网结果}

表\ref{tab:cigre}展示了CIGRE中压配电网上的实验结果。此处的协调挑战有所不同：系统管理的不是大型发电机组组合，而是数百个聚合为VPP的小型分布式能源。

\begin{table}[t]
\centering
\caption{CIGRE中压配电网实验结果。}
\label{tab:cigre}
\begin{tabular}{lcccc}
\toprule
\textbf{方法} & \textbf{成本} & \textbf{消纳率} & \textbf{电压} & \textbf{耗时} \\
 & (万元) & (\%) & 越限数 & (s) \\
\midrule
SCUC-MILP         & 0.92 & 85.4 & 3 & 18.7 \\
MADDPG-VPP        & 0.78 & 93.1 & 1 & 5.6 \\
\textbf{GridSynergy} & \textbf{0.71} & \textbf{96.8} & \textbf{0} & \textbf{6.3} \\
\bottomrule
\end{tabular}
\end{table}

GridSynergy在零电压越限的前提下实现了96.8\%的新能源消纳。在配电网层面，规划Agent的第四阶段交流潮流校验被证明至关重要：若缺少此步骤，SCUC-MILP调度将产生3个电压越限，因为MILP的直流潮流近似无法捕获高比例分布式能源接入配电网中的无功-电压耦合效应。

\section{消融实验}

为量化每个架构组件的贡献，我们在IEEE 30节点系统上通过选择性移除或替换GridSynergy组件进行消融实验。

\begin{table}[t]
\centering
\caption{IEEE 30节点系统消融实验结果。}
\label{tab:ablation}
\begin{tabular}{lccccc}
\toprule
\textbf{变体} & \textbf{成本} & \textbf{消纳} & \textbf{N-1} & \textbf{切负荷} & \textbf{耗时} \\
 & (万元) & (\%) & (\%) & (MWh) & (s) \\
\midrule
完整GridSynergy    & 12.85 & 97.3 & 100.0 & 0.08  & 9.1  \\
w/o 博弈Agent      & 13.42 & 93.5 & 100.0 & 1.82  & 7.2  \\
w/o 验证Agent      & 12.78 & 97.8 & 82.1  & 0.12  & 6.8  \\
w/o 分析Agent      & 12.90 & 97.1 & 100.0 & 2.15  & 14.5 \\
w/o LLM引擎        & 13.18 & 95.4 & 100.0 & 0.45  & 8.9  \\
纯LLM（无优化）    & 14.55 & 88.3 & 94.2  & 2.08  & 38.7 \\
纯优化（无LLM）    & 13.30 & 94.8 & 100.0 & 0.85  & 11.2 \\
\bottomrule
\end{tabular}
\end{table}

\textbf{无博弈Agent：}将博弈Agent替换为简单比例分配规则后，切负荷量增至1.82 MWh（恶化22.75倍），成本上升4.4\%。这确认了智能多VPP协调在故障场景下的关键作用。

\textbf{无验证Agent：}移除沙箱N-1检查后，N-1通过率降至82.1\%，意味着近五分之一的故障若按计划执行将导致安全越限。该结果凸显了验证Agent在运行安全中的不可或缺性。

\textbf{无分析Agent：}缺少分析Agent的异常检测能力时，故障检测延迟从5秒增至8.2秒（仅依赖简单阈值越限），切负荷增至2.15 MWh，因为延迟响应使频率偏差进一步扩大。

\textbf{无LLM引擎：}移除LLM引擎后，操作员需手动在结构化格式下指定全部约束。虽然数学结果仍保持良好（100\% N-1通过率、较低成本），但系统失去了自然语言接口和生成可读解释的能力。

\textbf{纯LLM与纯优化对比：}纯LLM变体（Grid-Agent风格）产生了最差的经济和安全表现。纯优化变体具备基本能力，但缺乏处理以自然语言表述的新型约束类型的灵活性。混合双引擎设计实现了二者最优组合。

\section{讨论}

\subsection{关键发现}

实验揭示了若干重要洞察：

\textbf{混合优于单一：}LLM推理与数学优化的结合始终优于二者单独使用。LLM提供灵活性和可解释性，优化引擎提供严谨性和保证。牺牲任何一方都会导致显著的性能退化。

\textbf{安全需要显式验证：}即使先进的AI系统也可能生成违反N-1准则的调度计划。验证Agent的沙箱化、穷尽式故障扫描对运行安全至关重要，不可被学习的启发式规则替代。

\textbf{多智能体协同产生协同效应：}每个Agent处理调度问题的不同方面，其协同——而非竞争——实现了全面覆盖。规划Agent提供初始计划，验证Agent确保其安全性，分析Agent检测偏差，博弈Agent协调响应——共同形成完整的感知-规划-验证-执行闭环。

\textbf{速度对实时运行至关重要：}GridSynergy完整流水线在10秒内完成，使其可用于实时调度。这一速度得益于并行沙箱N-1扫描、高效的MATD3推理（执行时仅需前馈网络推断）以及将重型优化委托给商业级求解器的设计。

\subsection{局限性}

当前系统的若干局限性值得讨论：

\textbf{向更大电网的泛化：}虽然GridSynergy在30节点和CIGRE网络上表现优异，但其在数千节点系统（如波兰2383节点系统或全规模省级电网）上的可扩展性仍需进一步验证。主要瓶颈是SCUC的MILP求解器运行时间随系统规模呈超线性增长。

\textbf{LLM可靠性：}尽管使用温度0.1，LLM引擎偶尔仍会误判复杂约束描述。我们通过结构化中间表示和验证检查加以缓解，但对于安全关键型应用中的LLM系统，完全可靠性仍是一个开放挑战。

\textbf{训练数据依赖性：}MATD3策略基于历史/合成场景离线训练。超出训练分布的极端事件可能导致次优博弈结果。带有安全约束的在线微调是一个有前景的方向。

\textbf{网络安全考量：}自然语言接口引入了新的攻击面：对抗性提示词可能操纵LLM引擎生成不安全的调度方案。我们建议采用纵深防御策略，包括输入净化、输出约束强制执行（验证Agent提供硬安全屏障）和对高影响决策的人工审批。

\subsection{未来工作}

\begin{enumerate}
    \item \textbf{分层MARL：}将博弈Agent扩展为支持大规模电网的分层协调架构（省级 $\to$ 区域级 $\to$ VPP级）。
    \item \textbf{在线学习：}通过安全在线强化学习使博弈Agent和分析Agent持续适应变化的电网工况。
    \item \textbf{多模态输入：}引入气象预报、卫星图像和市场电价等额外输入模态，实现更全面的决策。
    \item \textbf{联邦部署：}探索联邦学习架构，使VPP能在不共享敏感运行数据的前提下协同训练策略。
    \item \textbf{形式化验证：}应用形式化方法验证整体Agent协同协议的安全属性。
\end{enumerate}

\section{结论}

本文提出了GridSynergy——一个面向新能源电网的自主调度多智能体协同决策系统。系统的关键创新包括：（1）LLM推理与数学优化协同融合的混合双引擎架构；（2）四大专业协同Agent（规划、验证、博弈、分析）覆盖完整调度生命周期；（3）确保运行安全的沙箱化N-1安全验证；（4）基于MATD3的多VPP博弈协商实现近最优协调；（5）降低操作门槛的自然语言交互。

在IEEE 30节点系统和CIGRE配电网上的实验结果表明，GridSynergy相比传统SCUC方法降低发电成本12.0\%，实现97.3\%的可再生能源消纳率，在故障场景下维持近乎零负荷切除（0.08 MWh）。消融实验确认了每个架构组件——博弈Agent、验证Agent、分析Agent和LLM引擎——对系统整体性能的不可或缺贡献。

GridSynergy向完全自主的智能电网调度系统迈出了重要一步，能够在保持现代社会所需可靠性标准的前提下，充分发挥可再生能源的潜力。随着人工智能与电力系统持续融合，我们相信多智能体协同框架将在未来智能电网的运行中发挥日益核心的作用。

\section*{致谢}

作者感谢匿名审稿人提出的建设性意见。本研究得到了国家自然科学基金和国家电网公司科技项目的部分资助。

\begin{thebibliography}{00}

\bibitem{ire2025}
J. Okafor, D. Mensah, 等, ``Artificial Intelligence Agents for Smart Grid Applications: A Comprehensive Survey,''
\emph{IRE Journals}, 卷~9, 期~1, 页码~87--104, 2025.

\bibitem{llm_suc_2025}
H. Zhang, J. Wang, 等, ``LLM-SUC: Large Language Models for Security-Constrained Unit Commitment in Power Systems,''
\emph{arXiv preprint arXiv:2502.10557}, 2025.

\bibitem{marl_grid_2026}
R. Patel, M. Kumar, 等, ``Multi-Agent Reinforcement Learning for Power Grid Resilience under High Renewable Penetration,''
\emph{PLoS One}, 卷~21, 期~3, 页码~e0298541, 2026.

\bibitem{madrl_vpp_2026}
T. Rodriguez, A. Mueller, 等, ``Multi-Agent Deep Reinforcement Learning for Virtual Power Plant Scheduling in Distribution Networks,''
\emph{Sustainability}, 卷~18, 期~2, 页码~452--470, 2026.

\bibitem{topology_marl_2026}
K. Nakamura, L. Zhang, 等, ``Topology-Aware Multi-Agent Reinforcement Learning for Power Grid Operation,''
\emph{International Journal of Low-Carbon Technologies (IJLCT)}, 卷~21, 页码~234--248, 2026.

\bibitem{grid_agent_2025}
S. Li, Y. Chen, 等, ``Grid-Agent: A Large Language Model Agent Framework for Power Grid-Interactive Tasks,''
\emph{arXiv preprint arXiv:2508.05702}, 2025.

\bibitem{maddpg}
R. Lowe, Y. Wu, A. Tamar, 等, ``Multi-Agent Actor-Critic for Mixed Cooperative-Competitive Environments,''
见 \emph{Advances in Neural Information Processing Systems (NeurIPS)}, 2017, 页码~6379--6390.

\bibitem{matd3}
J. Ackermann, V. Gabler, T. Osa, 与 M. Sugiyama, ``Reducing Overestimation Bias in Multi-Agent Domains Using Double Centralized Critics,''
\emph{arXiv preprint arXiv:1910.01465}, 2019.

\bibitem{td3}
S. Fujimoto, H. van Hoof, 与 D. Meger, ``Addressing Function Approximation Error in Actor-Critic Methods,''
见 \emph{Proc. International Conference on Machine Learning (ICML)}, 2018, 页码~1587--1596.

\bibitem{wood_power}
A.~J.~Wood, B.~F.~Wollenberg, 与 G.~B.~Sheblé, \emph{Power Generation, Operation, and Control}, 第3版. Wiley, 2013.

\bibitem{stochastic_scuc}
Q. Wang, J. Wang, 与 Y. Guan, ``Stochastic Unit Commitment with Uncertain Demand Response,''
\emph{IEEE Trans. Power Systems}, 卷~28, 期~1, 页码~562--563, 2013.

\bibitem{cigre_benchmark}
K. Strunz, E. Abbasi, 与 R. Fletcher, ``Benchmark Systems for Network Integration of Renewable and Distributed Energy Resources,''
CIGRE Task Force C6.04.02, 技术报告 575, 2014.

\bibitem{langchain}
LangChain 文档. [在线]. 可访问: \url{https://python.langchain.com/}

\bibitem{pandapower}
L. Thurner, A. Scheidler, 等, ``pandapower---An Open-Source Python Tool for Convenient Modeling, Analysis, and Optimization of Electric Power Systems,''
\emph{IEEE Trans. Power Systems}, 卷~33, 期~6, 页码~6510--6521, 2018.

\bibitem{gurobi}
Gurobi Optimization, LLC, ``Gurobi Optimizer Reference Manual,'' 2024. [在线]. 可访问: \url{https://www.gurobi.com/}

\end{thebibliography}

\end{document}
